\section{The analytic completion of a one-weight Bessel propagation}
\label{sec:bessel-weight-propagation-v50}

This section proves the original final reconstruction step
\texttt{FinStatement}, the period-differentiation statement
\texttt{PerStatement}, and the generated-module Capelli eigenvalue statement
\texttt{CapValStatement}.  These are substantive producers used in both
\texttt{StepDown} and \texttt{StepUp}.  The full one-weight propagation is
still conditional on the complementary-minor evaluation and shifted
lowest-type quotient constructions.  We state that boundary explicitly
at the end.  None of the arguments below replaces the original integral
by an abstract branching assertion.

We retain BNV's notation $n=2r+k+1$, $G_H=U(H_{pk})$,
$G_\epsilon=U(J_\epsilon)$, $H_z$, $N$, and the actual character
$\psi_{z,c}$.  Put
\[
 \rho(g)=\|g\|_{\rm F},\qquad
 q_0=\min(p,q),\qquad
 a=\frac{h}{8\sqrt{2q_0}},\qquad \beta=a-1.
 \tag{BSP.1}
\]
The division in this notation is used only after positive rank has been
proved.  A Cayley matrix satisfies $P^*J_\epsilon P=H_{pk}$.
For a scalar function $F$ on matrices and a vector $\xi\in V_\pi$, the
original period is
\[
 \mathcal P_F(x)=
 \int_{H_z\times N}
 F\bigl(P(h'n'x)P^{-1}\bigr)\,
 \overline{\psi_{z,c}(n')}\,\pi(h')^{-1}\xi\,
 d\mu_W(h')d\mu_N(n').
 \tag{BSP.2}
\]
This is the source's \texttt{per}; the Haar measures retain their original
weak predicates, treated by AMD8.  Every statement concerning Bessel data
below allows every original $z,c,u,\pi,\mu_W,\mu_N,P$, with
$z$ a negative middle vector, all $c_a\ne0$, and the exact
$\texttt{IsHCTypeW}(k,r,z,u,\pi)$.

The inherited inputs are DC1--DC13, including the proved kernel Capelli
identity and degree \emph{orthogonality}, and BNV1.  No formula for the
formal degree, no Bessel multiplicity-one theorem, and no one-weight
propagation theorem is assumed.  Source evidence is pinned in the
accompanying manifest.

% BEGIN BSP1
\paragraph{BSP1. The whole original \texttt{GapNeedStatement}.}
\label{claim:BSP1}
For every $k,r\ge0$ and signature with $p=k+r$, $q=r+1$, the original
$\texttt{StepGap}(\epsilon,r,h)$ implies
\[
 q_0\ge1,\qquad h>h_{\rm need}(\epsilon),\qquad
 \beta>4\max(2nq_0,2rn^2)>0.
 \tag{BSP.3}
\]
In particular it supplies DC's large-gap construction threshold and the
integrability inequalities $\beta/2>2nq_0$, $\beta/4>2rn^2$.

\begin{proof}
The two hypotheses are exactly
\[
 \max(3n/2,1+npq)<h-2,\qquad
 4\max(2nq_0,2rn^2)<h/(8\sqrt{2q_0})-1.
\]
If $q_0=0$, the second inequality, with the source's zero-division
convention, would say that a nonnegative number is less than $-1$.
Thus $q_0\ge1$.  Put $s=8\sqrt{2q_0}\ge8$ and
$X=\max(2nq_0,2rn^2)\ge2nq_0\ge2$.
The second hypothesis is $h>s(4X+1)$.
It supplies the third member of $h_{\rm need}$ in (BNV.11), since
$q-1=r$.  Moreover
\[
 s(4X+1)\ge64nq_0+s>16nq_0+s+2,
\]
which supplies the second member.  The first hypothesis supplies the
first member directly.  This proves $h>h_{\rm need}$ and (BSP.3).
Notice that $k=0$ was not removed from the source's scope: its impossible
rank-zero case is excluded by the hypotheses themselves.
\end{proof}
% END BSP1

% BEGIN BSP2
\paragraph{BSP2. The exact kernel datum and rigidity of its cyclic lowest type.}
\label{claim:BSP2}
Here the dimension $n$ and signature $\epsilon$ are arbitrary; the
signature restrictions of BSP1 are not imposed.
Let $\lambda$ be Harish--Chandra, with noncompact gaps at least $h$,
$q_0\ge1$, and $h>h_{\rm need}(\epsilon)$.
For every original compact type $T$ of weight
$\Lambda=\lambda+\rho_n-\rho_c$ and every $\mu=\lambda+\rho$, there
exist a complete Hilbert space $\mathcal H$, a unitary representation
$\sigma$, an invariant smooth derived domain $D$, matrix-unit operators
$E_{ab}$ on $D$, and a compact irreducible subspace $V\subset D$ of
highest weight $\Lambda$, together with an isomorphism $\iota:V\to T$,
such that
\[
 \begin{split}
 C_Gv&=(\|\lambda\|^2-\|\rho\|^2)v\quad(v\in V),\\
 \langle v',\sigma(g)v\rangle
   &=\langle\iota v',\Psi_{T,\mu}(g)\iota v\rangle
       \quad(g\in G_\epsilon,\ v,v'\in V).
 \end{split}
 \tag{BSP.4}
\]
The subspace $V$ is stable under the actual compact group, not merely its
Lie algebra.  This is the entire original \texttt{PsiDatumStatement}.

More generally, every smooth compact irreducible $V$ of this highest
weight and Casimir in a unitary representation generates a cyclic closed
Hilbert subspace unitarily equivalent to the constructed model.
Its lowest type corresponds to $T$, and its coefficients satisfy the
four-vector orthogonality relation with some $d>0$ for the same Haar
measure.  This latter assertion is used only at the displayed
large-gap scope.

\begin{proof}
For existence, take DC's constructed Hilbert model and its smooth-vector
domain.  Let $j:T\to\mathcal H$ be the lowest-type embedding of DC7.
Divide it by the square root of the positive scalar in DC11, so that
it is an isometry and its coefficients are exactly $\Psi_{T,\mu}$.
Let $V=j(T)$ and $\iota=j^{-1}$.  DC3--DC4 give the derived action and
Casimir, DC7 gives smoothness, and the compact equivariance gives actual
compact-group stability.  These are all fields of the original datum.

Here is the rigidity argument for the stated broader input; in
particular we do not assume that this input is already discrete series.
DC8 applies to any such unitary derived representation: the fixed
Casimir and gap kill the negative first noncompact tensor layer.
DC9 then puts every noninitial compact type in strictly higher
coercive PBW layers, and makes the original type occur once.
The energy estimate and convergent orbit Taylor series in DC10 apply to
the same finite initial compact type and scalar Casimir.  They show that
the closure of the generated Lie module is group invariant and retains
that unique compact type.  If a closed invariant subspace is projected
orthogonally onto, its projection commutes with the compact group.
Its restriction to the unique initial type is zero or identity by
Schur's lemma; cyclicity then makes the whole projection zero or
identity.  Thus the cyclic closure is irreducible.

Identify the initial type isometrically with the initial type of DC's
model and take its diagonal copy in the direct sum of the two unitary
representations.  It has the same weight and Casimir.  The preceding
PBW, closure, and irreducibility argument applies to this diagonal copy.
Its two coordinate projections are bounded nonzero intertwiners with
dense images onto the two cyclic closures.  Composing either projection
with its adjoint gives a positive scalar times the identity, by
irreducibility; the scalar is nonzero.  Hence each projection is a
nonzero scalar multiple of a unitary isomorphism.  The prescribed
isometry on the initial types makes the two scalars agree.  We obtain
a unitary equivalence identifying those initial types and intertwining
the actual group actions.
Finally DC12 transports its full coefficient orthogonality and positive
degree to this cyclic subspace.  This uses the proved Hilbert-space
orthogonality, not an unreconstructed degree polynomial.
\end{proof}
% END BSP2

% BEGIN BSP3
\paragraph{BSP3. The whole original \texttt{CapValStatement}.}
\label{claim:BSP3}
For arbitrary $n,\epsilon,\lambda,h,T,\mu$ satisfying BSP2's
Harish--Chandra, noncompact-gap, $h>h_{\rm need}$, $q_0\ge1$,
and weight hypotheses, let $(\mathcal H,\sigma,D,E,V,\iota)$ be
\emph{any} complete unitary derived datum with compact irreducible
type of weight $\lambda+\rho_n-\rho_c$, its stated Casimir,
actual compact-group stability, and the exact coefficient identity
(BSP.4).  The ambient representation need not be irreducible or be
the one constructed in BSP2.
Let $S=U(\mathfrak{gl}_n)V\subset D$ be the generated module, and put
$\theta_i=\lambda_i-(n-1)/2$.  Then for every $\omega\in\mathbb C$
and every $m\in S$,
\[
 \mathcal C_E(\omega)m
       =\prod_{i=0}^{n-1}(\omega+\theta_i)\,m.
 \tag{BSP.5}
\]
The Capelli operator uses the original ordered column determinant
\[
 \mathcal C_E(\omega)
 =\sum_{\tau\in S_n}\operatorname{sgn}(\tau)
   \prod_{j=0}^{n-1}
       \bigl(E_{\tau(j),j}+\delta_{\tau(j),j}(\omega-j)\bigr),
 \tag{BSP.6}
\]
with the rightmost operator acting first.  Thus in particular
$\mathcal C_E(-\theta_s)$ kills the entire generated module, as needed
by both one-weight steps.

\begin{proof}
The hypotheses include the exact coefficient identity (BSP.4);
$D$ is an invariant derived domain, so every finite derivative word
may be applied to its vectors.  Differentiating that identity on the
right and using DC4 gives, for $v,v'\in V$,
\[
 \langle v',\sigma(g)\mathcal C_E(\omega)v\rangle
   =p_\lambda(\omega)\langle v',\sigma(g)v\rangle,\qquad
 p_\lambda(\omega)=\prod_i(\omega+\theta_i).
 \tag{BSP.7}
\]
The shift is exact: substitute $Y=\omega-(n-1)/2$ in DC4's
$\mathcal C(Y+(n-1)/2)$ identity.

Set $u=\mathcal C_E(\omega)v-p_\lambda(\omega)v$.
Let $\mathcal K$ be the closed span of the group translates of $V$.
It contains every $E_{ab}m$ for $m\in D\cap\mathcal K$:
real infinitesimal actions are limits of difference quotients of group
orbits in the closed subspace, and complex matrix units are complex
linear combinations of real Lie algebra elements.  Induction shows
$S\subset\mathcal K$, so $u\in\mathcal K$.
Equation (BSP.7) and unitarity make $u$ orthogonal to every group
translate of $V$; hence $u=0$.  This proves (BSP.5) on $V$ without
assuming that the ambient representation is irreducible.

The ordered determinant (BSP.6) is central in the enveloping algebra.
This is the centrality calculation included in DC4: commuting a matrix
unit through the ordered column expansion pairs the off-diagonal
commutator terms by the corresponding transposition, with the
$(\omega-j)$ shifts canceling the diagonal commutators.  Consequently
$\mathcal C_E(\omega)$ commutes with every $E_{ab}$ in every
representation of the matrix-unit relations.  Therefore the kernel of
$\mathcal C_E(\omega)-p_\lambda(\omega)$ is invariant under all the
$E_{ab}$ and contains $V$.  By the definition of the generated module
it contains $S$, proving the full assertion.  Evaluation at
$\omega=-\theta_s$ makes the $s$th factor zero.
\end{proof}
% END BSP3

% BEGIN BSP4
\paragraph{BSP4. The original invariant standard and dual multipliers.}
\label{claim:BSP4}
For every middle vector $z$, without requiring negativity, set
$x_* = e_0$ if $r>0$, and $x_*=z$ if $r=0$.
For a column $y$ and a row $\alpha$ define
\[
 m_y(g)=x_*^*H_{pk}gy,\qquad
 m^\vee_\alpha(g)=\alpha g^{-1}x_*.
 \tag{BSP.8}
\]
These are exactly the original \texttt{cMul} and \texttt{cMulD}.
For every $s\in S_z=H_zN$ and $g\in G_H$ they satisfy
\[
 m_y(sg)=m_y(g),\qquad m^\vee_\alpha(sg)=m^\vee_\alpha(g).
 \tag{BSP.9}
\]
For the complexified right derivative $R_X$,
\[
 R_Xm_y=m_{Xy},\qquad
 R_Xm^\vee_\alpha=-m^\vee_{\alpha X}.
 \tag{BSP.10}
\]
For every $y,\alpha$ there is one constant $C$ such that both absolute
values in (BSP.8) are at most $C\sum_{i,j}|g_{ij}|$ on $G_H$.
These are all conclusions of the original \texttt{MultStatement}.

\begin{proof}
If $r>0$, the defining isotropic flag makes every element of $N$ fix
its first vector $e_0$, and $H_z$ fixes the entire outer frame.
If $r=0$, $N=\{1\}$ and $H_z$ fixes $z$.  Thus $sx_*=x_*$ in either
case.  Since $s$ is $H_{pk}$-unitary, also
$x_*^*H_{pk}s=x_*^*H_{pk}$.  This proves both identities (BSP.9).
Differentiating $g\exp(tX)y$ gives $gXy$, while differentiating
$(g\exp(tX))^{-1}x_*=\exp(-tX)g^{-1}x_*$ gives
$-Xg^{-1}x_*$.  This proves (BSP.10) for real Lie directions and then
for their complexification; equivalently both expressions have local
holomorphic extensions on invertible matrices, so the same derivatives
are the source's ambient complex derivatives.
The first norm bound is a finite sum of fixed coefficients times entries
of $g$.  On $G_H$ one has $g^{-1}=H_{pk}g^*H_{pk}$, and $H_{pk}$ is
a signed permutation matrix.  Hence the sum of absolute entries of
$g^{-1}$ equals that of $g$, giving the second bound.  Take the larger
of the two constants.  The argument includes $z=0$ and all rank-zero
cases allowed in the original multiplier statement.
\end{proof}
% END BSP4

% BEGIN BSP5
\paragraph{BSP5. The whole original \texttt{PerStatement}.}
\label{claim:BSP5}
Assume the original signature identities and \texttt{StepGap}.
If $F$ is continuous on $G_\epsilon$ and
$|F(g)|\le C\rho(g)^{-\beta}$ there, then (BSP.2) is absolutely
integrable for every original datum, every $x\in G_H$, and every $\xi$.
If $F$ is smooth in the source's local ambient sense and every
complexified right derivative word $E_{\mathbf w}F$ satisfies
\[
 |E_{\mathbf w}F(g)|\le C_{\mathbf w}\rho(g)^{-\beta},
 \tag{BSP.11}
\]
then differentiation commutes with the actual period:
\[
 \mathcal P_{E_{ab}F}(x)
    =\left.
       E_{ab}\bigl(y\longmapsto\mathcal P_F(P^{-1}yP)\bigr)
      \right|_{y=PxP^{-1}}.
 \tag{BSP.12}
\]
The operator on the right is the original \texttt{eabB}.

\begin{proof}
BSP1 gives $\beta>2rn^2$.  Conjugation by $P$ changes Frobenius norm
by fixed positive comparison constants.  BNV1 therefore makes
$\rho(P h'n'xP^{-1})^{-\beta}$ integrable.  The remaining character
and unitary factors in the period have norm $\|\xi\|$.
Continuity on the supported subgroups gives strong measurability,
and proves the first assertion.

For a real Lie direction $X$ in $G_\epsilon$, write
$x_t=xP^{-1}\exp(tX)P$.  The scalar integrand is then
$F(P h'n'xP^{-1}\exp(tX))$.
For $|t|\le t_0$ the inverse of $\exp(tX)$ has bounded Frobenius norm,
so uniformly in $h',n'$,
\[
 \rho(P h'n'xP^{-1}\exp(tX))
       \ge c\,\rho(P h'n'xP^{-1}),\qquad c>0.
\]
The derivative of the scalar integrand is
$(R_XF)(P h'n'xP^{-1}\exp(tX))$ and has the same integrable majorant,
by (BSP.11) and the finite expression of $X$ in the complex matrix units.
The fundamental theorem of calculus bounds its difference quotients
by that majorant.  Dominated convergence differentiates (BSP.2) under
the integral.  It also gives continuity of the resulting derivative.
Repeating this argument with every derivative word proves smoothness
in group coordinates and commutation with every real derivative.
Complex linearity yields (BSP.12) with exactly the source's matrix-unit
convention.

For completeness these conclusions have the local ambient meaning
required on the right side of (BSP.12).  Near any group point,
$G_\epsilon$ is a smooth embedded real submanifold of the invertible
matrices.  A local smooth retraction onto it extends any smooth function
on the group to an ambient smooth function.  Complexified tangential
derivatives of this extension are independent of the retraction.
Thus the calculated derivatives are the literal source derivatives,
not derivatives of an unrelated representative.  Haar support
eliminates all off-group values of the original functions.
\end{proof}
% END BSP5

% BEGIN BSP6
\paragraph{BSP6. The group power-integral threshold used by the final Fubini step.}
\label{claim:BSP6}
For any signature with $q_0\ge1$ and every $b>2nq_0$,
\[
             \int_{G_\epsilon}\rho(g)^{-b}\,dg<\infty.
 \tag{BSP.13}
\]
The Haar normalization is arbitrary.  On $G_\epsilon$ also
$\rho(g)\ge1$, $\rho(g^{-1})=\rho(g)$, and
$\rho(gh)\le\rho(g)\rho(h)$.

\begin{proof}
The last three assertions follow from Frobenius submultiplicativity,
$g^{-1}=J_\epsilon g^*J_\epsilon$, and $gg^{-1}=I$.
For the integral use the polar coordinates of DC6:
$g=\kappa e^X$ with $\kappa$ compact and
$X=X^*$, $J_\epsilon XJ_\epsilon=-X$.
Put $s=\|X\|_{\rm F}$.  Its nonzero eigenvalues are
$t_1,-t_1,\ldots,t_{q_0},-t_{q_0}$, allowing zeros, so
$\max_i|t_i|\ge s/\sqrt{2q_0}$ and
$\rho(e^X)\ge\exp(s/\sqrt{2q_0})$.
The polar Jacobian already calculated in DC6 is at most
$\exp(n\sqrt{2q_0}s)$.
Its remaining radial shell factor is polynomial in $s$, since the
noncompact tangent space is finite dimensional.  The resulting
integrand is bounded at infinity by a polynomial times
\[
 \exp\left[-\frac{b-2nq_0}{\sqrt{2q_0}}s\right],
\]
which is integrable.  Near zero the polar Jacobian is continuous.
The compact factor has finite Haar mass.  This proves (BSP.13),
including the exact threshold used in the source's
\texttt{finB\_intJ}.  A positive rescaling of Haar measure changes only
the finite constant.
\end{proof}
% END BSP6

% BEGIN BSP7
\paragraph{BSP7. Pointwise reproduction inside the actual right-regular space.}
\label{claim:BSP7}
Let $\mathcal K$ be a closed invariant subspace of $L^2(G_\epsilon)$
whose coefficients have the four-vector orthogonality identity with
degree $d>0$.  Suppose $F$ is continuous on $G_\epsilon$, is in $L^2$,
and its class $f$ belongs to $\mathcal K$.
For every $x\in G_\epsilon$,
\[
 \int_{G_\epsilon}F(y)
        \langle f,R(y^{-1}x)f\rangle\,dy
           =\frac{\|f\|_2^2}{d}F(x).
 \tag{BSP.14}
\]
This is an equality of continuous functions, not just of $L^2$ classes.

\begin{proof}
Haar measure on $G_\epsilon$ is invariant under right translation and
inversion.  Indeed the polar decomposition reduces the modulus of the
adjoint determinant to the compact factor, where it is one, and to
$\exp X$, where it is $\exp(\operatorname{tr}(\operatorname{ad}X))=1$:
the matrix-unit roots of the complexified Lie algebra occur in opposite
pairs, so the trace is zero.  Thus the group is unimodular.
Write $\gamma(g)=\langle f,R(g)f\rangle$.
Orthogonality makes $\gamma\in L^2$, so the integral in (BSP.14)
is absolutely convergent by Cauchy--Schwarz.  It is continuous in $x$
by strong continuity of translation in $L^2$.
Let $\varphi$ be any continuous compactly supported test function.
For fixed $x$, substitute $y=zt$ and apply Fubini:
\[
 \begin{split}
 &\int \overline{\varphi(z)}
       \left(\int F(y)\gamma(y^{-1}zx)\,dy\right)dz\\
 &\qquad=\int
       \langle\varphi,R(t)f\rangle\,
       \overline{\langle R(x)f,R(t)f\rangle}\,dt.
 \end{split}
 \tag{BSP.15}
\]
The double absolute integral is at most
$\|\varphi\|_1\|F\|_2\|\gamma\|_2$, which justifies this interchange.
If $\Pi$ is the orthogonal projection onto $\mathcal K$, replace
$\varphi$ by $\Pi\varphi$ in the first coefficient.  This is allowed
because all $R(t)f$ lie in $\mathcal K$.
Four-vector orthogonality makes the last line of (BSP.15)
\[
 \frac{\|f\|_2^2}{d}\langle\Pi\varphi,R(x)f\rangle
       =\frac{\|f\|_2^2}{d}
          \int\overline{\varphi(z)}F(zx)\,dz.
\]
Thus the difference of the two continuous functions in (BSP.14),
translated by $x$, pairs to zero with every compactly supported
continuous test function.  If it were nonzero somewhere, a sufficiently
small neighborhood and a nonnegative test function there would give
a nonzero pairing (multiply first by the conjugate phase at that
point).  Haar measure is positive on nonempty open sets.  This
contradiction proves the pointwise identity.
\end{proof}
% END BSP7

% BEGIN BSP8
\paragraph{BSP8. The full original \texttt{finProj} reconstruction.}
\label{claim:BSP8}
Assume the signature identities, \texttt{StepGap}, and a Harish--Chandra
$\lambda'$ with noncompact gaps at least $h$.
Let $W_f$ be an original linear space of scalar matrix functions, and
$W'$ a compact irreducible subspace of the smooth vectors of
$L^2(G_\epsilon)$, with highest weight
$\Lambda'=\lambda'+\rho_n-\rho_c$ and Casimir $c_{\lambda'}$.
Assume actual compact-group stability, and the source's exact image
relation
\[
 v\in W'\ \Longleftrightarrow\
 \exists F\in W_f:\ v=[F]_{L^2}
 \quad\text{for every smooth }v.
 \tag{BSP.16}
\]
Every $F\in W_f$ is required to be continuous on the group and bounded
by $C_F\rho^{-\beta}$.
For every such $F$ whose class is a nonzero element of $W'$, there
exist an original compact type $T'$, $\mu'=\lambda'+\rho$, and
$v_1,v_2\in T'$ such that
\[
 F(x)=\int_{G_\epsilon}F(y)
      \langle v_2,\Psi_{T',\mu'}(y^{-1}x)v_1\rangle\,dy
       \qquad(x\in G_\epsilon).
 \tag{BSP.17}
\]

\begin{proof}
BSP1 supplies the large gap for BSP2.
The cyclic Hilbert space generated by $W'$ is therefore unitarily
equivalent to the constructed model, with the entire lowest type
matched.  In particular it has four-vector orthogonality for some
$d>0$, and there is $t\in T'$ with
\[
 \langle f,R(g)f\rangle
       =\langle t,\Psi_{T',\mu'}(g)t\rangle,\qquad f=[F]_{L^2}.
\]
Here the lowest-type embedding has been normalized isometrically,
as in BSP2.  The function $F$ is genuinely in $L^2$:
its displayed bound, BSP1, and BSP6 show this directly.
Consequently the formal \texttt{toL2} in (BSP.16) is its usual class,
and $f\ne0$.
Apply BSP7 and take
$v_2=t$, $v_1=(d/\|f\|_2^2)t$.
The scalar is nonzero and acts in the linear inner-product slot.
The result is (BSP.17), with no almost-everywhere ambiguity and with
the exact kernel required by the original \texttt{finProj}.
\end{proof}
% END BSP8

% BEGIN BSP9
\paragraph{BSP9. The full original final Fubini and nonvanishing extraction.}
\label{claim:BSP9}
Under the original signature, gap, Bessel-data, and kernel hypotheses,
suppose $F$ is continuous on $G_\epsilon$ and
$|F(y)|\le C\rho(y)^{-\beta}$.
For every $x\in G_H$ the following double integrand is Bochner integrable:
\[
 F(y)\,
 \langle v_2,\Psi_{T',\mu'}
        (y^{-1}P h'n'xP^{-1})v_1\rangle\,
 \overline{\psi(n')}\,\pi(h')^{-1}\xi.
 \tag{BSP.18}
\]
The measures are $dy\,d\mu_W(h')\,d\mu_N(n')$.
If in addition $F$ satisfies (BSP.17) and
$\mathcal P_F(x)\ne0$, then some $x_0\in G_H$ satisfies
\[
 A_{T',\mu',P}(v_1,v_2,\xi;x_0,x)\ne0
 \tag{BSP.19}
\]
for the original integral (BNV.5), with the same Bessel datum.
These are the original \texttt{finFub\_int'} and \texttt{finFub}
conclusions; BNV1 proves their \texttt{IntSStatement} input.

\begin{proof}
Put $g=P h'n'xP^{-1}\in G_\epsilon$ on the Haar support.
The Frobenius coefficient estimate proved in BNV19, or its proof
from DC6, gives
$|\langle v_2,\Psi(y^{-1}g)v_1\rangle|
       \le C'\rho(y^{-1}g)^{-a}$.
Since $\rho\ge1$ and $a=\beta+1$, this is at most
$C'\rho(y^{-1}g)^{-\beta}$.
Let $A=\rho(y)$, $B=\rho(y^{-1}g)$, and $G_0=\rho(g)$.
Then $A,B\ge1$ and $G_0\le AB$.
The following elementary majorant suffices:
\[
 A^{-\beta}B^{-\beta}
   \le G_0^{-\beta/4}A^{-\beta/2}.
 \tag{BSP.20}
\]
Indeed $G_0^{-\beta/4}\ge A^{-\beta/4}B^{-\beta/4}$,
and $A^{-\beta}B^{-\beta}\le
A^{-3\beta/4}B^{-\beta/4}$ because $A,B\ge1$.
This slightly simpler majorant replaces the source's sum of two
majorants without changing any required exponent.

The first factor on the right of (BSP.20) is integrable over
$H_z\times N$ by BNV1 and norm comparison under $P$:
$\beta/4>2rn^2$ by BSP1.
The second is integrable over $G_\epsilon$ by BSP6, since
$\beta/2>2nq_0$.  Their product is therefore integrable.
The character and representation factors have norm $\|\xi\|$;
the supported functions are continuous and strongly measurable.
This proves the full Bochner-integrability assertion (BSP.18).

Insert (BSP.17) into (BSP.2) and interchange the integrals using
that assertion.  With $x_0=P^{-1}yP$ one has exactly
\[
 y^{-1}P h'n'xP^{-1}
       =P(x_0^{-1}h'n'x)P^{-1}.
\]
Thus the inner integral after the interchange is
$A(v_1,v_2,\xi;P^{-1}yP,x)$.
If all these vector integrals were zero, the resulting expression
for $\mathcal P_F(x)$ would be zero, a contradiction.
Since $P^{-1}yP\in G_H$, this proves (BSP.19) with the original
two free translates and all original normalization factors.
\end{proof}
% END BSP9

% BEGIN BSP10
\paragraph{BSP10. The whole original \texttt{FinStatement}.}
\label{claim:BSP10}
Take every datum and hypothesis in BSP8, together with every original
Bessel datum and a fixed $\xi\in V_\pi$.  Assume the source's
period-injectivity condition
\[
 \begin{gathered}
 F\in W_f,\quad
 \mathcal P_F(x)=0\ \text{for every }x\in G_H\\
 \Longrightarrow\quad F(g)=0\ \text{for every }g\in G_\epsilon.
 \end{gathered}
 \tag{BSP.21}
\]
Then $\texttt{BesselNZ}(\epsilon,\lambda',u)$ holds.
This proves the full original \texttt{FinStatement}, including its
$\texttt{IntSStatement}$ antecedent by BNV1, every original space
$W_f,W'$, every original Haar pair, and the exact period (BSP.2).

\begin{proof}
The nonzero irreducible compact type $W'$ contains a nonzero vector
$f$.  The exact image relation (BSP.16) supplies $F\in W_f$ representing
it.  If all periods of $F$ vanished, (BSP.21) would make its function
zero on the group and hence its $L^2$ class zero, contradicting $f\ne0$.
Choose $x$ with $\mathcal P_F(x)\ne0$.
BSP8 supplies $T',\mu',v_1,v_2$ and the pointwise reproduction formula;
BSP9 supplies $x_0$ with the nonzero original vector integral.
These, together with the unchanged $P,z,c,\pi,\mu_W,\mu_N,\xi,x$,
are every field in the existential definition of
$\texttt{BesselNZ}(\epsilon,\lambda',u)$.
The construction has not assumed existence of such a period or of
a shifted discrete series beforehand: nonvanishing came from the
explicit injectivity hypothesis (BSP.21), and the shifted model came
from its stated lowest-type and Casimir data.
\end{proof}
% END BSP10

\subsection{What remains in each one-unit propagation}

The new complete producers are the original
\texttt{GapNeedStatement}, \texttt{PsiDatumStatement},
\texttt{CapValStatement}, \texttt{MultStatement},
\texttt{PerStatement}, \texttt{FinStatement}, and the intermediate
pointwise reconstruction and final Fubini theorems used to prove the last
one.  The source's two complete step proofs still require the following
substantive inputs.

First, the standard and dual tensor coefficient maps must be shown to be
smooth $L^2$ functions with the exact derivative bounds.  Their formulas
are fixed: with $x_1=P x_*$,
\[
 \begin{split}
 T_f(g)&=\sum_j (x_1^*J_\epsilon g e_j)
                  \langle w,\sigma(x_0^{-1}g)f_j\rangle,\\
 T_f^\vee(g)&=\sum_j (e_j^*g^{-1}x_1)
                  \langle w,\sigma(x_0^{-1}g)f_j\rangle.
 \end{split}
 \tag{BSP.22}
\]
Here $w$ lies in the original lowest compact type and every $f_j$ in
its generated module.  BSP4 verifies the exact scalar multipliers and
their derivatives; it does not by itself prove all derivative bounds
for these generated-module coefficients.

Second, the complementary Capelli minors must produce the relevant
Euler-operator eigenspaces in the standard and dual tensors.
BSP3 proves the required annihilation
$\mathcal C_E(-\theta_s)S=0$, but the noncommutative adjugate identities
turning it into those eigenspace vectors remain separate.
The original \texttt{EvalPStatement} and \texttt{EvalMStatement} then
assert a nonzero period of a specific such tensor image whenever
the old period is nonzero and
$u_j\ne\lambda_s+1/2$ or $u_j\ne\lambda_s-1/2$, respectively.
Their products of excluded factors alone are not a proof of this
evaluation assertion.

Third, \texttt{QuotPStatement} and \texttt{QuotMStatement} must produce
spaces $W_f,W'$ for the shifted weights $\lambda+e_s$ and $\lambda-e_s$,
with exactly (BSP.16), the shifted Casimir, actual compact stability,
the smooth decay properties, and the period-injectivity (BSP.21).
BSP5 supplies differentiation through their actual periods and BSP10
completes the last step once these spaces are constructed.  It does
not assert the quotient existence.

Consequently the whole original \texttt{StepStatement} remains open.
BNV23, the whole \texttt{BesselNVGap}, and the supplied-form global setting
remain conditional on that source producer.  No selected global
obligation is closed by this section.
